% ==================================================================
% id:       func-019
% topic:    פונקציות
% subtopic: תחום הגדרה ותכונות של פונקציות
% solution: full
% ==================================================================
\textbf{בדקו את זוגיות הפונקציות הבאות:}
\begin{enumerate}\bMs{2}
\item  $\ds{f(x)=\frac{x^2}{\sqrt{x^4+1}}}$
\begin{solbox}
ראשית נשים לב שהפונקציה מוגדרת לכל הישר הממשי. והיא זוגית כמנה של פונקציות זוגיות.
\end{solbox}
\item  $f(x)=x+1$
\begin{solbox}
לא זוגית ולא אי זוגית דוגמה נגדית:
$f(1)=2 , f(-1)=0$
\end{solbox}
\item  $f(x)=x^3$
\begin{solbox}
אי זוגית,
מאחר וחזקה אי זוגית היא פונקציה אי זוגית.
\end{solbox}
\item  $\ds{f(x)=\log_2\lr{\frac{1-x^3}{1+x^3}}}$
\hint{חשבו את $f(-x)$ והיעזרו בתכונות הלוגריתם.}
\begin{solbox}
הפונקציה מוגדרת כאשר
$\frac{1-x^3}{1+x^3}>0$,
כלומר עבור
$-1<x<1$,
וזהו תחום סימטרי סביב $0$.
לכל $x$ בתחום מתקיים
$$f(-x)=\log_2\lr{\frac{1+x^3}{1-x^3}}=\log_2\lr{\lr{\frac{1-x^3}{1+x^3}}^{-1}}=-\log_2\lr{\frac{1-x^3}{1+x^3}}=-f(x),$$
ולכן הפונקציה אי-זוגית.
\end{solbox}
\item $f(x)=\sin^2(x)+\cos(x)$
\hint{חשבו את $f(-x)$ והיעזרו בזוגיות או באי-הזוגיות של $\sin$ ושל $\cos$.}
\begin{solbox}
פונקציה זוגית,
כסכום של פונקציות זוגיות.
ידוע כי
$\cos(x)$
זוגית
ו
$sin(x)$
אי זוגית ולכן
$sin^2(x)$
זוגית
כי
\\
$sin^2(-x)=(sin(-x))^2=(-sin(x))^2=sin^2(x)$
\end{solbox}
\eMs
\end{enumerate}
